\documentclass[10pt,a4paper]{amsart}
\usepackage[margin=19mm]{geometry}
\usepackage{amsmath,amssymb,mathtools}
\usepackage[scr=rsfs,cal=euler]{mathalfa}
\usepackage{xfrac}
\usepackage[unicode,hidelinks,bookmarks=false]{hyperref}
\newcommand{\ie}{i.e.\ }
\newcommand{\cf}{cf.\ }
\newcommand{\resp}{respectively\ }
\newcommand{\Subsection}{\subsection}
\providecommand{\nobreakdash}{\nobreak\hbox{-}}
\setlength{\emergencystretch}{2em}
\multlinegap0pt
\usepackage{bm}
\usepackage{bbm}
\usepackage{dsfont}
\usepackage{xspace}
\usepackage{cancel}
\usepackage{mathtools}
\usepackage{amsthm}
\makeatletter
\@ifundefined{Th}{\newtheorem{Th}{Th}}{}
\@ifundefined{Proposition}{\newtheorem{Proposition}[Th]{Proposition}}{}
\makeatother
\DeclareMathOperator{\R}{\mathbb{R}}
\title[Constrained Variable Projection for Structured Problems]{Constrained Variable Projection for Structured Problems}
\author{Emanuele Zangrando, Sara Venturini, Francesco Rinaldi, Francesco Tudisco}
\date{Source: arXiv:2606.23939v1 (CC BY 4.0)}
\begin{document}
\maketitle

\noindent\emph{Source and licence.} Constrained Variable Projection for Structured Problems. Version arXiv:2606.23939v1 (CC BY 4.0), distributed under the Creative Commons Attribution 4.0 International licence, as recorded in the arXiv licence metadata for this exact version and shown on the abstract page https://arxiv.org/abs/2606.23939v1. Full rights notice: licenses/arxiv-2606.23939-CC-BY-4.0.txt.\par\smallskip

\noindent\textit{Editorial scope.} Counted unit: the Proposition whose source label is \texttt{prop:differentiability\_optimalconstrained} in the article (source0.tex lines 357--367, source label \texttt{prop:differentiability\_optimalconstrained}), delivered together with the article's own complete original proof of it (source0.tex lines 369--505, one \texttt{proof} environment of 137 lines). No mathematics is written, repaired or re-derived: statement and proof are the article's own text line for line. Required context retained uncounted: none. Every definition, display and label of those lines is retained unchanged. Omitted from this package and not redistributed: 1--356, 368--368, 506--1541. Nothing inside a retained excerpt is omitted or altered: every retained body is delivered exactly as the article's lines are written, so no omission receipt of any kind accompanies this package. Citations: the retained excerpts carry no citation command at all, so no bibliography entry of the article is transcribed and no reference is redistributed. Reference adaptation: none was needed and none was made; every reference inside the delivered unit resolves inside it, so no unresolved reference or citation is produced. Printed slips kept literal: no printed word, symbol, hypothesis or number of the counted statement or of its proof was altered. Layout and class: the article's own class is \texttt{article} with packages including arxiv, amsmath, amsfonts, amssymb, amsthm, lipsum, graphicx, epstopdf, mathrsfs, algorithmic, enumitem, algorithm, graphics, hyperref; the wrapper is the lane's pinned \texttt{amsart} editorial wrapper at 10pt on A4, which restates the theorem-like environments the retained text needs and adds the article's own macro definitions that the retained text needs (\texttt{\textbackslash R}), with extra wrapper packages (bm, bbm, dsfont, xspace, cancel, mathtools). The delivered numbering is the unit's own. Assets: no retained span contains a figure environment, a graphics-inclusion command, a PDF-page inclusion, a table or a TikZ picture; no image file is copied into this package, the package declares no asset dependency and no image file is redistributed with it. Licence: version arXiv:2606.23939v1 of this article is distributed under the Creative Commons Attribution 4.0 International licence, as recorded in the arXiv licence metadata for this exact version and shown on the abstract page https://arxiv.org/abs/2606.23939v1. A full rights notice is delivered at licenses/arxiv-2606.23939-CC-BY-4.0.txt. Characters: every retained excerpt is pure ASCII, so the delivered engine, XeLaTeX with its default Latin Modern text font, reports no missing character.
\begin{Proposition}(Differentiability of optimal solution for constrained problems)\label{prop:differentiability_optimalconstrained} 
Let $\mathscr W$ be a Chebyshev subset of $E:=\R^N$ and consider the problem
    \[
    \mathcal P_{\mathscr W}^{Q}(\widehat w):= \arg\min_{w \in \mathscr W} \|w-\widehat w \|_{Q}^2.
    \]
Then, for $Q = (M^\top M + \lambda \Omega^\top  \Omega), \| x\|_Q^2 = \|Q^{1/2}x \|_2^2, \widehat w = Q^{-1} M^\top y$, and assuming that $\lambda > 0, m \geq N$ and $\Omega$ full rank, the map 
\[
\Pi : \R^{m \times N} \times \R^{m} \to \R^{N},\quad  \Pi(M,y) = \mathcal P^{Q(M)}(\widehat w(M,y))
\]
is differentiable for all $(M,y) \in \R^{m \times N} \times \R^{m}$ if and only if $\mathscr W$ is affine.
\end{Proposition}

\begin{proof}
Let \(E=\mathbb R^{N}\), endowed with the $\ell^2$ inner product.
Since \(\mathscr W\) is a Chebyshev subset of a finite-dimensional Euclidean
space, \(\mathscr W\) is closed and convex. Hence, for every positive definite
matrix \(Q\), the \(Q\)-metric projection onto \(\mathscr W\) is well-defined
and single-valued.

Because \(\Omega ^\top\Omega\succ0\) and \(\lambda>0\), we have
\[
Q(M)=M^\top M+\lambda\Omega^\top\Omega\succ0
\]
for every \(M\). Thus, \(Q(M)^{-1}\) depends smoothly on \(M\), and so does $
\widehat w(M,y)= Q(M)^{-1}M^\top y.$

We first prove the easy direction. Suppose that \(\mathscr W\) is affine, i.e., $\mathscr W=w_*+\mathscr L,$ where \(\mathscr L\subset E\) is a linear subspace. Fix a basis
\(v_1,\dots,v_r\) of \(\mathscr L\). For fixed \(Q\succ0\), the projection of
\(\widehat w\) onto \(\mathscr W\) has the form
\[
\mathcal P_{\mathscr W}^Q(\widehat w)
=
w_*+\sum_{i=1}^r \alpha_i v_i.
\]
The coefficients are determined by the normal equations
\[
\langle
\Big(\widehat w-w_*-\sum_{j=1}^m \alpha_j v_j\Big)
,Q v_i
\rangle=0,
\qquad i=1,\dots,r.
\]
Equivalently, $G(Q)\alpha=b(Q,\widehat w),$ where
\[
G(Q)_{ij} = \langle v_i,v_j \rangle_{Q},
\qquad
b(Q,\widehat w)_i=\langle\widehat w-w_*, Q v_i \rangle.
\]
Since \(Q\succ0\), the Gram matrix \(G(Q)\) is positive definite, hence
invertible. Therefore
\[
\alpha=G(Q)^{-1}b(Q,\widehat W)
\]
depends smoothly on \((Q,\widehat w)\). Since \(Q(M)\) and
\(\widehat w(M,y)\) depend smoothly on \((M,y)\), the map
\[
\Pi(M,y)=\mathcal P_{\mathscr W}^{Q(M)}(\widehat w(M,y))
\]
is differentiable, indeed smooth.

Conversely, assume that \(\Pi\) is differentiable for every \((M,y)\).
Choose \(M_0\in\mathbb R^{m\times N}\) with \(\operatorname{rank}(M_0)=N\),
which is possible because \(N \leq m\). Set $Q_0:=Q(M_0)\succ0$ and consider the linear map
\[
L:\mathbb R^{m}\to \mathbb R^{N},
\qquad
L(y)= Q_0^{-1} M_0^\top y.
\]
Since \(\operatorname{rank}(M_0)=N\), the matrix \(Q_0^{-1} M_0^\top\) has rank
\(N\), hence \(L\) is surjective. Thus \(L\) admits a linear right inverse
\(R:\R^N\to\mathbb R^{m}\).
For fixed \(M_0\), the map $
y\mapsto \Pi(M_0,y) = \mathcal P_{\mathscr W}^{Q_0}(L(y))$
is differentiable by assumption. Since \(L\circ R=\operatorname{Id}_E\), we get $\mathcal P_{\mathscr W}^{Q_0}(z) =
\Pi(M_0,Rz)$.
Therefore, the full \(Q_0\)-metric projection $
z\mapsto \mathcal P_{\mathscr W}^{Q_0}(z)$ is differentiable everywhere on \(E\).

Now reduce to the ordinary Frobenius projection. Define the invertible linear
map
\[
T:E\to E,\qquad T(w)=Q_0^{1/2}w,
\]
and set $C:=T(\mathscr W)=Q_0^{1/2}\mathscr W .$
Since 
$\|w-\widehat w\|_{Q_0} =
\|T(w)-T(\widehat w)\|_2,$
we have $
T\bigl(\mathcal P_{\mathscr W}^{Q_0}(\widehat W)\bigr) =
\mathcal P_C(T\widehat W),$
or equivalently
\[
\mathcal P_C
=
T\circ \mathcal P_{\mathscr W}^{Q_0}\circ T^{-1}.
\]
Hence the ordinary Frobenius projection \(\mathcal P_C\) is differentiable
everywhere.

We now show that this forces \(C\) to be affine. Let \(c\in C\), and let $
\mathcal K_cC:=\overline{\operatorname{cone}(C-c)}
$
be the tangent cone of \(C\) at \(c\). For closed convex sets in a finite
dimensional Hilbert space, the directional derivative of the metric projection
at a point \(c\in C\) is given by $d\mathcal P_C(c;h)=\mathcal P_{\mathcal K_cC}(h),$
where \(\mathcal P_{\mathcal K_cC}\) denotes the orthogonal projection onto the closed
convex cone \(\mathcal K_cC\).

But \(\mathcal P_C\) is Frechet differentiable at \(c\), so
$h\mapsto d\mathcal P_C(c;h)$
is linear, i.e., \(h\mapsto \mathcal P_{\mathcal K_cC}(h)\) is linear. The range of
this linear map is exactly \(\mathcal K_cC\), which is therefore a linear subspace.
Let now $A:=\operatorname{aff}(C)$
be the affine hull of \(C\), and let $
V:=A-c=\operatorname{span}(C-c)
$
be its associated direction space. Since \(\mathcal K_cC\subset V\) and
\(C-c\subset \mathcal K_cC\), we get
$
V=\operatorname{span}(C-c)\subset \mathcal K_cC\subset V.
$
Therefore, $
\mathcal K_cC=V
$
for every \(c\in C\).

We claim that every \(c\in C\) is in the relative interior of \(C\) inside
\(A\). Indeed, if some \(c\in C\) were not in \(\mathrm{int}_A(C)\), the
supporting hyperplane theorem would give a nonzero \(u\in V\) such that $\langle u,z-c\rangle_2\le 0\  
  \text{for every }z\in C
$. Passing to the tangent cone gives
$
\langle u,h\rangle_2\le 0\,\,
 \text{for every }h\in \mathcal K_cC.
$
But \(\mathcal K_cC=V\), and \(u\in V\), so choosing \(h=u\) gives $
\|u\|_2^2\le0,
$
a contradiction. Hence \(C=\mathrm{int}_A(C)\).
Thus \(C\) is both relatively open and relatively closed in its affine hull
\(A\). Since \(A\) is connected and \(C\neq\emptyset\), it follows that $
C=A.
$
Therefore \(C\) is affine. Finally, since \(T^{-1}\) is an invertible linear map, $
\mathscr W=T^{-1}(C)
$
is affine as well.
This proves the converse direction, and hence the proposition.
\end{proof}

\end{document}
